% SPDX-License-Identifier: CC-BY-SA-4.0
% Author: Matthieu Perrin

\begingroup

\begin{exercice}[Pile bornée]
  On cherche à implémenter une pile bornée comme suit : 
  
  \begin{lstlisting}
    public class BoundedStack<T> {
      private final AtomicReferenceArray<T> array;
      public BoundedStack(int capacity) {
        array = new AtomicReferenceArray<>(capacity);
      }
      public void push(T value) {
        for(int i = 0; i<array.length(); i++) {
          if(array.compareAndSet(i, null, value)) return;
        }
      }
      public T pop() {
        for(int i = array.length()-1; i>=0; i--) {
          T value = array.get(i);
          if(value != null && array.compareAndSet(i, value, null)) 
          return value;
        }
        return null;
      }
    }
  \end{lstlisting}

  \begin{question}
  \item Décrire l'exécution de l'algorithme, quand un processus appelle, séquentiellement :\\ \lstinline{p.push(a); p.push(b); p.pop(); p.push(c);}\\
    sur une pile $p$ de capacité 3. 
  \item L'algorithme est-il linéarisable ? 
  \item Quelles sont les conditions de progression vérifiées par l'algorithme ? 
  \end{question}

\end{exercice}

\endgroup
\endinput
